\scanpage{145}
\diagram{$(**)\quad$\begin{tikzcd}
1 \arrow[r,no head] & V'_K \arrow[r,no head] \arrow[d,leftrightarrow] & S_K^0 \arrow[r,no head] \arrow[d,leftrightarrow,"\unsure{i_{K/K'}|i_{K,K'}}"] & G_m \arrow[r,no head] \arrow[d,leftrightarrow,"\lambda\mapsto\lambda^r"] & 1 \\
1 \arrow[r,no head] & V'_{K'} \arrow[r,no head] & S_{K'}^0 \arrow[r,no head] & G_m \arrow[r] & 1
\end{tikzcd}\qquad $r=[K':K]$}

On trouve:

\uline{Corollaire 3} Le système projectif des $S_K^0$, pour $K$ sous-extension
finie de $\overline{\mathbb Q}$, est can. isom (via norme,
au syst. projectif des $S_K^0$ pour $K\subset\overline{\mathbb Q}$, $K$ weilien
(et galoisien si on veut), ou encore au syst. projectif des
\deleted{$\unsure{V_K}$} $T_K/V_{K_0\deleted{\unsure{R}}}$, où
$K\subset\overline{\mathbb Q}$ est Weilien et $K_0=K\cap\mathbb R$.

\uline{Corollaire 4} Soit $K$ ext. finie de $\mathbb Q$. Il existe un et un seul hom
\[
G_m\xrightarrow{i}S_K^0
\]
tel que $\varepsilon i(\lambda)=\lambda^2$.

Comme $\operatorname{Ker}\varepsilon$ est compact, tout hom
$G_m-\operatorname{Ker}\varepsilon$ est nul, donc l'unicité est claire.
Pour l'existence, on peut supposer grâce à la prop. que $K$ est weilien.
Si $K\subset\mathbb R$ \inserted{i.e. $K_0=K$}, on a
$S_K^0\xrightarrow[\sim]{\varepsilon}G_m$, et il suffit de prendre $i(\lambda)=\lambda^2$.
Si $K_0\ne K$, alors on prend pour
\[
i=i_{K_0,K}:S_{K_0}^0=G_m-S_K^0,
\]
en utilisant la commutativité de $(**)$ (où $r=2$).
\margin{144}
